\section{Definition}
\label{sec:definition}

Let $(A,B)$ have variances $\sigma_A^2$ and $\sigma_B^2$, correlation $\rho$, and variance ratio $\kappa=\sigma_B^2/\sigma_A^2$.
Write $D=A-B$.

\begin{definition}[Paired Efficiency Factor]
\label{def:pef}
\[
  \eta
  = \frac{\sigma_A^2+\sigma_B^2}{\Var(D)}
  = \frac{1+\kappa}{1+\kappa-2\sqrt{\kappa}\,\rho}.
\]
\end{definition}

The quantity is a ratio of second moments.
It does not require a Gaussian margin.
Equivalently,
\begin{equation}
  \eta = \frac{1}{1-u},
  \qquad
  u = \frac{2\sigma_{AB}}{\sigma_A^2+\sigma_B^2},
  \label{eq:pef_u}
\end{equation}
so $u$ is the pairwise concordance correlation of Lin with the location-shift term set to zero~\parencite{lin1989}, and also the intraclass correlation for the pair.
The following properties are algebraic consequences of~\cref{def:pef}.

\begin{enumerate}[label=P\arabic*.]
  \item $\eta>1$ if and only if $\rho>0$, for every $\kappa>0$.
        The boundary $\eta=1$ is the line $\rho=0$.
  \item $\eta(\kappa,\rho)=\eta(1/\kappa,\rho)$.
  \item $\eta\in[\tfrac12,\infty)$.
        The minimum $\tfrac12$ is attained at $\rho=-1$ and $\kappa=1$.
        In general $\eta(\kappa,1)=(1+\kappa)/(1-\sqrt{\kappa})^2$ and $\eta(\kappa,-1)=(1+\kappa)/(1+\sqrt{\kappa})^2$.
  \item For fixed $\rho>0$, $\eta$ is maximised at $\kappa=1$, where $\eta=1/(1-\rho)$.
        For fixed $\rho<0$, $\eta$ is minimised at $\kappa=1$.
  \item $\eta\to 1$ as $\kappa\to 0$ or $\kappa\to\infty$.
  \item $\eta=1/(1-u)$ with $u$ as in~\cref{eq:pef_u}.
\end{enumerate}

These properties decide the efficiency of the paired difference as an estimator of a mean difference.
They do not decide whether $D$ is a linearly sufficient predictor of a binary outcome.
